\documentclass[11pt,a4paper]{ctexart}

% -------- 宏包加载 --------
\usepackage{amsmath, amssymb, amsfonts} % 数学公式与符号
\usepackage{graphicx}                  % 图片插入
\usepackage{grffile}
\usepackage{geometry}                  % 页面边距设置
\usepackage{booktabs}                  % 三线表
\usepackage{caption}                   % 图表标题控制
\usepackage{subcaption}                % 子图支持
\usepackage{float}                     % 强制固定图表位置 [H]
\usepackage{cite}                      % 参考文献引用
\usepackage{titlesec}                  % 标题格式定制
\usepackage{bm}                        % 数学符号加粗
\usepackage{siunitx}                   % 国际单位制格式化

% -------- 页面与布局设置 --------
\geometry{top=25mm, bottom=25mm, left=30mm, right=30mm}

% -------- 超链接宏包 --------
\usepackage{hyperref}                  
\hypersetup{
    colorlinks=true,
    linkcolor=blue,
    citecolor=blue,
    urlcolor=magenta
}



% -------- 标题样式修改 --------
\titleformat{\section}{\large\bfseries\color{black}}{\thesection}{1em}{}
\titleformat{\subsection}{\boldmath\bfseries}{\thesubsection}{1em}{}

% -------- 正文开始 --------
\begin{document}

% -------- 论文题目 --------
\title{\Large\bfseries 基于菲涅尔波带片的水波聚焦波场设计研究报告}
\author{潮汐工坊}
\date{ }
\maketitle


% -------- 摘要与关键词 --------
\begin{abstract}
菲涅尔波带片是一类利用波的干涉与衍射实现能量聚焦的重要波前调控器件，其理论已广泛应用于光学领域，而在宏观水波体系中的研究仍相对有限。本文将菲涅尔波带片理论推广至水表面波体系，结合 MATLAB 数值计算完成波带片参数优化。实验中搭建了机械驱动直线波源、水槽及基于快速棋盘格解调（Fast Checkerboard Demodulation，FCD）的光学测量系统，实现了水面高度场的非接触式重建，并测定了实验条件下的水波衰减参数。在此基础上，设计并制作了振幅型、相位 X 型、相位 h 型以及带缓冲区的正弦平方式相位 h 型四种菲涅尔波带片，对其聚焦特性进行了实验测量与 MATLAB 仿真对比，并采用二维高斯拟合和聚焦效率等指标进行了定量评价。结果表明，实验获得的聚焦图案与理论预测及数值模拟基本一致，验证了菲涅尔波带片在水波体系中实现波前调控与能量聚焦的可行性。其中，相位型波带片整体聚焦性能优于振幅型，h 型结构兼顾较高聚焦效率与较好的参数鲁棒性；引入边界缓冲区后，虽然聚焦效率提升有限，但能够有效减弱边缘散射与背景扰动，提高聚焦波场质量。本研究为水波波前调控及其在粒子操控、波浪能汇聚和流体波场设计等方向的进一步应用提供了实验基础和参考。

\vspace{0.8em}
\noindent\textbf{关键词：}菲涅尔波带片；水波聚焦；波前调控；波带片设计；聚焦效果
\end{abstract}
\newpage
\tableofcontents
\newpage

% -------- 一、引言 --------
\section{引言}

\subsection{历史研究}

菲涅尔波带片（Fresnel zone plate, FZP）是一类基于波的干涉与衍射原理实现聚焦的人工结构。与传统透镜依靠介质折射率改变传播方向不同，波带片通过对传播路径上的波场振幅或相位进行空间调制，使不同位置传播到焦点处的波产生相长干涉，从而实现类似透镜的聚焦效果。由于其聚焦机制来源于波动本质，因此波带片并不局限于可见光领域，而能够推广至声波、水波、电磁波以及其他形式的机械波中。

菲涅尔波带片的发展根源于十九世纪初的波动光学研究。1818年，法国物理学家 Fresnel 在研究光的衍射现象时提出了基于惠更斯原理的半波带理论，将复杂的衍射问题转化为不同空间区域波的相位叠加问题\cite{Fresnel1818}。根据该理论，来自同一波阵面的不同区域到达空间某一点时具有不同的传播距离，其相位差决定了最终光强分布。若人为选择性地保留或阻挡满足特定相位关系的区域，则可以增强目标位置处的干涉强度，这一思想构成了菲涅尔波带片的理论基础。

随后，波带片理论逐渐发展成为经典傅里叶光学的重要组成部分。Born 和 Wolf 在《Principles of Optics》中系统总结了菲涅尔衍射理论、半波带理论以及波带片聚焦条件，指出波带片本质上是一种利用空间周期结构实现波前调控的衍射元件\cite{BornWolf1999}。与传统折射透镜相比，波带片避免了材料折射率和曲率对聚焦能力的限制，其焦距主要由波长以及结构尺寸决定，因此在短波长领域具有特殊优势。

随着现代光学的发展，Goodman 从傅里叶光学角度进一步阐述了波带片的作用机制，将其视为一种空间频率滤波器。通过对入射波场进行周期性的振幅或相位调制，波带片能够改变光场的空间频谱分布，并在传播过程中形成特定的衍射重构图样\cite{Goodman2005}。这一观点使波带片从传统衍射器件发展成为现代波前调控技术的重要组成部分。

二十世纪以来，随着微纳加工技术的发展，菲涅尔波带片逐渐由理论研究进入实际应用阶段。Firester 等人研究了微型菲涅尔波带片的制造方法以及聚焦性能，实验验证了通过周期性透明区域调制光场可以实现高质量聚焦\cite{Firester1973}。此后，波带片被广泛应用于极紫外光学、X 射线显微成像以及高分辨率光学系统中。例如，在短波长成像领域，由于传统透镜材料受到吸收和色散限制，波带片成为实现纳米尺度聚焦的重要光学元件\cite{Boivin1977}。

进一步的研究表明，菲涅尔波带片的物理本质并不依赖于光波本身，而是依赖于所有波动系统中普遍存在的干涉和衍射规律。只要传播介质满足线性波动方程，并且能够对波场进行空间调制，就可以构造类似的波带片结构。因此，菲涅尔波带片的概念逐渐从光学领域推广至声学、水波以及弹性波等研究方向。

水波作为一种典型的宏观机械波体系，具有传播速度较低、波长较大、实验可视化程度高等特点，非常适合作为研究波动现象的平台。相比于光学实验中需要复杂探测设备才能观察传播过程，水面波能够直接通过波纹、干涉条纹以及驻波结构体现波场变化，因此成为研究衍射、干涉和波前调控的重要实验体系。早期水波实验主要集中于验证经典波动理论，例如圆孔衍射、双缝干涉以及驻波模式等，而近年来随着人工结构设计的发展，水波逐渐成为研究人工波场调控的新平台。

从这一发展历程来看，菲涅尔波带片经历了由光学理论提出、微结构器件实现，再到多种波体系推广的发展过程。其核心思想始终保持不变：通过人为设计空间结构，使传播过程中的不同波前分量在目标位置满足相位匹配条件，从而实现能量汇聚。本课题利用水表面波作为实验对象，将这一经典光学思想推广至宏观流体体系，通过设计不同结构参数的水波菲涅尔波带片，实现水面波场的聚焦，为进一步研究水波操控以及微粒输运等问题提供基础。

\subsection{研究现状}

近年来，随着人工结构设计以及超材料（metamaterials）理论的发展，利用人工结构主动调控水波传播过程逐渐成为波动物理领域的重要研究方向。传统水波研究主要关注自然边界条件下波的传播、反射、折射以及散射等现象，而现代研究则更加关注通过设计具有特定几何结构的人工界面，实现对水波振幅、相位以及传播方向的主动控制。这一研究思路与光学中的波前调控思想高度相似，使水波成为研究人工光学类比系统的重要平台。

在水波人工调控领域，早期研究主要借鉴电磁超材料和声学超材料的设计思想，通过周期性结构改变水波传播过程中的等效参数，实现异常折射、负折射以及波束控制等现象。Farhat 等人利用人工结构实现了水面波的特殊传播控制，证明了通过合理设计边界条件可以改变表面波传播行为，并进一步实现水波隐身等功能\cite{Farhat2009}。该类工作表明，水波并非只能被动传播，而可以通过人工结构实现类似于光学器件中的主动调控。

随后，水波超材料和人工透镜方向得到快速发展。Zhang 等人利用零折射率人工结构实现了水波的宽带聚焦和准直传播，通过调节结构参数改变水波等效传播特性，实现了传统介质中难以获得的波场分布\cite{Zhang2014}。该研究展示了利用人工结构设计水波传播路径的可能性，也证明了宏观水波体系能够实现类似光学透镜的功能。

除利用等效介质参数调控水波外，基于相位调控的波前工程方法近年来受到广泛关注。在光学领域，超表面（metasurface）通过亚波长结构单元控制入射波相位，从而实现聚焦、偏转以及复杂波束生成等功能。类似思想被逐渐引入水波体系，通过设计不同高度、形状以及周期分布的结构单元，使水波在传播过程中获得预定相位延迟，从而形成特定空间分布的波场。该类研究说明，水波体系同样可以利用结构化设计实现类似光学元件的功能\cite{Che2024STVP}。

与上述基于连续介质参数调控的方法不同，菲涅尔波带片提供了一种更加直接的空间干涉调控方案。波带片并不需要改变传播介质本身，而是通过选择性增强或削弱特定空间位置的波贡献，使不同路径传播的波在焦点处满足相位匹配条件。因此，波带片具有结构简单、易于加工以及实验可观察性强等优势。近年来，类似思想已被推广到声学和水动力学体系中，用于实现机械波聚焦和能量汇聚。

在水波研究中，利用人工结构实现聚焦主要包括两类方法：一类是通过连续折射率分布或复杂曲面结构模拟传统透镜；另一类则是通过离散结构单元调控波的传播相位。其中，后者与菲涅尔波带片思想更加接近。通过设计不同区域的透射、反射或阻挡特性，可以改变不同空间位置水波的贡献，使焦点位置形成明显的能量增强区域。这种离散化结构不仅降低了制造难度，同时也便于研究结构参数（例如波带宽度、间距以及反射系数）对聚焦性能的影响。

近年来，随着对水波场精细调控需求的增加，研究者进一步探索了结构化水波与物质相互作用之间的关系。类似光学捕获和声学镊子的思想，研究人员开始利用水面波产生局域压力分布、稳定势阱以及定向输运效应，实现对漂浮微粒或颗粒体系的操控。例如，基于拓扑水波结构的研究表明，通过设计具有特殊相位分布的水波场，可以形成稳定的空间结构，并进一步影响颗粒运动轨迹\cite{Wang2024}。此外，预印本工作中动态水波镊子等研究进一步展示了利用时变水波场实现颗粒捕获、移动和操纵的可能性\cite{Wang2026}。

然而，与光学领域成熟的菲涅尔波带片研究相比，水波体系中的波带片研究仍处于发展阶段。目前已有研究更多集中于水波超材料、人工边界以及复杂结构调控，而针对经典菲涅尔波带片结构在水面波中的系统研究仍相对有限。特别是对于振幅型波带片，不同波带宽度、间距以及反射特性如何影响水波聚焦效果，以及单一线源激励条件下波场传播规律，仍具有进一步研究价值。

本课题基于上述研究背景，将经典菲涅尔波带片思想引入水面波体系，采用单一直线正弦型波源产生近似平面传播水波，并设计制作多种结构参数不同的水波菲涅尔波带片。实验中通过观察不同波带片对水面波传播的调制效果，验证其聚焦现象，并比较不同结构设计对聚焦位置和聚焦强度的影响。该研究不仅验证了菲涅尔波带片在宏观流体波体系中的适用性，同时为进一步利用结构化水波实现粒子操控、能量汇聚以及复杂波场构造提供了实验基础。

\subsection{文章结构}
本文主要分为八个部分，第二部分先介绍菲涅耳原理，重点讲解如何从光波类比到水波，两者的相似性与差异性，并指出水波研究相较于光波的困难；第三部分介绍实验装置，首先介绍观测装置，即水波平台，并简要概述快速棋盘格解调（FCD）原理，接着介绍设计的直线波源，最后讲述4种波带片的设计和迭代过程；第四部分介绍实验与仿真的结果，首先测量了水波的衰减系数，作为仿真的参数，接着给出4种波带片的实验结果与仿真结果，对比说明波带片的聚焦效果；第五部分对FCD程序测量的精度与不确定度进行说明；第六部分分析了实验的误差来源，主要分为程序上和实验上两个版块；第七部分对全文进行总结；第八部分提出未来计划进行的实验。此外，在附录中有对实验用到的所有仿真程序，以及使用的测量程序的结构说明。

% -------- 二、从光学到水波 --------
\section{从光学到水波}

\subsection{菲涅尔原理与波带片设计}

根据惠更斯-菲涅尔原理（Huygens-Fresnel Principle），空间中任一点的波动场是波前上各次级波源产生的次波相干叠加的结果。从点波源 S 发出的波在到达焦点 P 时，由于传播的光程（或波程）不同，会产生相位差异，进而发生干涉现象。

\begin{figure}[htbp]
    \centering
    \includegraphics[width=0.5\textwidth]{bodaipianshiyitu.png}
    \caption{振幅型菲涅尔波带片原理示意图}
    \label{fig:bodaipianshiyitu}
\end{figure}

为了在 P 处实现最大限度的相长干涉，从而实现能量叠加，可以将球面波波前划分为多个相位差为 $\pi$ 的环形波带，即半波带。当其次波到达焦点时，其波程差满足以下条件：
\begin{equation}
(\rho_m - R_m) - (\rho_0 - R_0) = \frac{m\lambda}{2}, \quad m=1, 2, 3, \dots
\end{equation}

在此条件下，主焦距 $f$ 近似定义为：
\begin{equation}
f = \frac{r_m^2}{m\lambda}
\end{equation}

根据对半波带调控方式的不同，波带片主要分为两种类型。第一种为振幅型波带片（如图~\ref{fig:bodaipianshiyitu} 所示），其通过交替阻挡偶数（或奇数）半波带，消除相位相反的相消干涉项，使到达焦点的所有次波均处于同相叠加状态，从而实现能量聚集。第二种为相位型波带片，该方案利用波在不同介质或不同水深中传播速度的差异，对相邻半波带引入 $\pi$ 的相位延迟。这种方式不阻挡波浪能量，而是通过台阶状或连续的微结构改变局部相速度，使原本会相消的波带发生相位反转，从而实现比振幅型更高的聚焦效率\cite{hecht2017optics}。

\subsection{波动的相似性与差异性}

在宏观波动力学中，光波与水波均遵循基本的波动方程，具备干涉、衍射等经典现象，且都可以通过改变空间相位差来调控波阵面。因此，我们可以将光学中成熟的菲涅尔波带片理论直接引入流体力学，通过在水下设计特定的几何结构来调控水波相位，实现对水波能量的聚焦。这为高效收集波浪能、设计新型海岸防护结构等工程应用提供了理论基础\cite{li2018concentrators}。

然而，在微观调控机制与物理边界上，水波与光波存在显著差异，主要体现在以下三个方面。首先是色散关系与相位调控。光波的相位调控通常依赖折射率的变化，而水波的相速度则强烈依赖于局部水深，水越浅，波速越慢。水波的完整色散关系需同时考虑重力与表面张力的影响，表示为：
\begin{equation}
\omega^2 = \left(gk + \frac{\sigma k^3}{\rho}\right)\tanh(kh_0)
\label{dispersion}
\end{equation}
其中，$\sigma$ 为表面张力系数，$\rho$ 为液体密度。在实验设计中，需要根据该方程逆向求解出波速随水深的变化规律，从而精确反演出台阶的高度与波带片的物理尺寸。其半波带半径 $r_n$ 需严格遵循：
\begin{equation}
r_n = \sqrt{n\lambda f + \left(\frac{n\lambda}{2}\right)^2}
\end{equation}
其次是能量衰减。水波除了色散关系更复杂外，还受到流体粘滞摩擦带来的线性振幅衰减影响。对于理想势流中的重力波，其阻尼系数 $\gamma$ 近似为：
\begin{equation}
\gamma = \frac{2\nu\omega^4}{g^2}
\end{equation}
其中，$\nu$ 为运动粘度。
最后是非线性效应。在流体力学中，水波流经锐利的地形边缘时极易发生流动分离并产生涡流，这会造成额外的能量耗散与杂散波干扰。

\subsection{水波菲涅尔波带片设计}

将理论模型转化为实际的水波透镜时，我们不可避免地会遇到真实的流体力学问题。我们的设计主要旨在解决以下三个技术难点。

第一是水位突变引发的反射问题。若仅满足理论上的相位要求，水位会呈现从深水区到极浅水区的急剧变化。这种实体台阶凸起的边缘会直接阻挡入射波，造成不可忽略的水波反射，导致大量能量在穿透波带片前就已损耗。

第二是涡流耗散。流体流经阶跃的锐利边缘时，会发生边界层流动分离并产生细小涡流。这种非线性效应不仅消耗水波动能，还会引发杂乱的二次辐射，破坏干涉场的纯净度。

第三是浅水区能量耗散与深水判据。为了积累足够的相位差，水波必须途经较浅的水域。而真实水波在浅水区受底部粘滞摩擦影响极大，线性衰减显著增加，因此应通过调整波带片构型保证其不在浅水区。针对本实验设定的 $6\text{Hz}$ 波源频率，代入水波色散关系方程计算可得其理论波长 $\lambda \approx 4.87\text{cm}$。根据经典水波理论的浅水波判据（通常界定为水深 $h < \lambda/20$），当进入极浅水区域时，非线性效应与底部摩擦耗散将急剧放大。因此，我们在反演地形高度时，通过计算强制设定了浅水区的安全底线（$3\text{mm}$），从而在确保相邻波带精确积累 $\pi$ 相位差的同时，有效避免了因水层过薄导致的波阵面崩溃与过度耗散\cite{li2021surface}。

在整体方案与参数匹配方面，本项目基于波动衍射理论，利用 MATLAB 编写了参数化设计程序，并结合 3D 打印技术制造了不同功能的波带片。针对 $6\text{Hz}$ 频率的水波，我们通过数值模拟优化了波带片级数 $n$ 与焦距 $f$ 的匹配关系，并重点开发了两类器件。第一类为振幅型波带片，其通过物理挡板阻挡偶数（或奇数）半波带实现能量叠加，但这不可避免地会阻挡约 50\% 的入射能量。第二类为相位型波带片，该方案运用水深变化改变局部相速度的原理，通过程序精密计算出实体结构，为相邻波带引入 $\pi$ 的相位差，成功实现了比振幅型更卓越的聚焦效果。为了进一步探究空间形貌对相位分布的调控作用，我们制作了两类特殊的相位型波带片：一类是固定波带片的高度、连续改变其长度 $x$，称为 $x$ 相位型波带片（如图~\ref{fig:sub_image_a} 所示）；另一类是固定波带片的长度、连续改变其高度 $h$，称为 $h$ 相位型波带片（如图~\ref{fig:sub_image_b} 所示）。

\begin{figure}[htbp]
    \centering
    \begin{subfigure}[b]{0.45\textwidth}
        \centering
        \includegraphics[width=\textwidth]{xxiangwei.png}
        \caption{$x$ 相位型波带片示意图}
        \label{fig:sub_image_a}
    \end{subfigure}
    \hfill 
    \begin{subfigure}[b]{0.45\textwidth}
        \centering
        \includegraphics[width=\textwidth]{hxiangwei.png}
        \caption{$h$ 相位型波带片示意图}
        \label{fig:sub_image_b}
    \end{subfigure}
    \caption{两类不同形貌的相位型波带片结构图}
    \label{fig:phase_zone_plates}
\end{figure}

% -------- 三、实验装置 --------
\section{实验装置}
\subsection{实验装置与测量系统}

本实验的物理装置具体结构如图 \ref{fig:实验装置} 所示。整个系统主要由水槽主体、造波模块以及光学观测模块三部分组成。所需核心器材包含：$90\text{cm} \times 60\text{cm} \times 10\text{cm}$（壁厚 $1\text{cm}$）的透明亚克力水槽、工业相机、标准棋盘格、LED 面光源、亚克力三棱柱、驱动电机及若干铝型材支架。

\begin{figure}[htbp]
    \centering
    \includegraphics[width=0.4\textwidth]{shiyanzhuangzhi.jpg}
    
    \caption{实验装置图}
    \label{fig:实验装置} 
\end{figure}

在激振造波方面，系统的波源机构由固定在铝型材支架上的电机与亚克力三棱柱构成。通过电机驱动浸入水中的亚克力三棱柱作周期性的竖直振动，以此在水槽内部激发频率稳定、波阵面平直的连续平面波。

在数据采集方面，水面波动的形态捕捉依托于底部的透射式光学系统。水槽正下方平铺放置一张标准棋盘格标定板（单个黑白网格的物理尺寸精确至 $5\text{mm} \times 5\text{mm}$），其下方紧贴放置一块与棋盘格尺寸等大的 LED 面光源，提供均匀的照明，配合架设在水槽上方的工业相机进行俯视拍摄，从而实现对水面波纹演化的成像与记录。

为了确保波源激发的波纹在达到稳定状态后再引入波带片，同时为水波的衍射与聚焦提供更充裕的空间演化范围与观测视野，我们将实验水槽的有效观测尺寸定为 $60\text{cm} \times 90\text{cm}$。广阔的水面区域不仅有效延缓了边界反射波到达中心观测区的时间，也尽可能避免了波源近场效应对波带片聚焦特性的干扰。

此外，背景波场的纯净度直接决定了聚焦效率的测量精度。由于水波的干涉对相位极其敏感，微弱的边界反射驻波即可与主干涉场发生叠加，从而干扰波带片本身的聚焦效果。前期初步实验的数据处理结果表明，仅在水槽垂直边缘敷设 $1\text{cm}$ 厚的海绵，其能量耗散能力有限，无法起到理想的消波作用。

为此，我们查阅相关文献，对水槽的边界条件进行了重新设计与优化\cite{Konispoliatis2021}。首先，在沿波浪传播方向（行波方向）的末端，我们借鉴了海岸工程中防波堤的消波原理，设置了厚度为 $5\text{cm}$、倾角为 $\arctan\left(\frac{30}{50}\right)$ 的斜坡海绵。当水波向斜坡传播时，由于水深逐渐变浅，波的相速度减小、波幅增大，最终发生变形甚至破碎。这一过程将波动的机械能大量转化为流体内部的粘滞耗散能，从而实现了高效的吸收与抑制端面反射。其次，对于水槽的两侧侧壁，我们将其直接替换为表面光滑的亚克力板。相比于粗糙的海绵，光滑的亚克力材质能最大程度地降低流体在侧壁边界层内的粘滞摩擦阻力。在宏观流体力学模型中，这种处理可近似视为理想的“滑移边界”或“自由边界”，从而保证了水波在沿主轴传播过程中的平直性与波前完整性。经过上述水槽尺寸的扩大与边界条件的优化，实验环境内的反射现象得到了显著抑制，为后续精确观察与评估菲涅尔波带片的聚焦性能提供了稳定、理想的物理平台。
% ==========================================

观测程序是参照Wildeman, Sander的文章\cite{Wildeman2018}，采用快速棋盘格解调（Fast Checkerboard Demodulation，FCD）技术，利用二维周期性图案--棋盘格作为背景，类似于频率调制技术，背景图案将作为高频载波，其相位由水面的变化调制，在明确规定的条件下，可采用傅里叶解调技术直接提取该信号。具体计算过程如下（程序介绍放在附录中）：

由于每个像素均由两个独立的分量组成，所以图案一定含有两个线性无关的载波矢量。类比晶体学，我们将一般的二维周期性图案写为调和级数的形式：
\begin{equation}
I_0(\bm{r})
= \sum_{m=-\infty}^{\infty}
  \sum_{n=-\infty}^{\infty}
  a_{mn}\,
  e^{i\,(m\bm{k}_1 + n\bm{k}_2)\cdot\bm{r}}
\end{equation}
其中，$\bm{k}_1 \times \bm{k}_2 \neq 0$，$\bm{k}_1$和$\bm{k}_2$称为倒易晶格向量，并且由于$I_0$是实数，$a(m, n) = a^*(-m, -n)$。


将（2）代入（1）得，
\begin{equation}
I(\bm{r})
= \sum_{m,n}
  a_{mn}\,
  e^{i\,(m\bm{k}_1 + n\bm{k}_2)\cdot(\bm{r}-\bm{u}(\bm{r}))}
\end{equation}
从（3）可以看出，位移场$\bm{u(r)}$的作用是调控每个谐波的相位。

如果失真信号的幅度不大，且构成该信号的波长也不过短，那么$\bm{u(r)}$在空间傅里叶域（k空间）中的影响将集中在载波峰值$\bm{k}_c \in n\bm{k}_1 +m\bm{k}_2$附近，接着采用傅里叶域滤波技术，从公式(2)和(3)中的求和项中分离出如下形式的复信号：
\begin{align}
g_0(\bm{r}) & = a_c\, e^{i\bm{k}_c\cdot\bm{r}} \\
g(\bm{r}) & = a_c\, e^{i\bm{k}_c\cdot(\bm{r}-\bm{u}(\bm{r}))}
\end{align}

则$\bm{u(r)}$可以通过下式得到， 
\begin{align}
\phi(\bm{r})
& \equiv \operatorname{Im}\!\bigl(
  \ln\bigl(g(\bm{r})\,g_0^*(\bm{r})\bigr)
\bigr)
= -\bm{k}_c\cdot\bm{u}(\bm{r}) \\
\phi_1(\mathbf{r}) &= -\mathbf{k}_1\cdot\mathbf{u}(\mathbf{r}) \\
\phi_2(\mathbf{r}) &= -\mathbf{k}_2\cdot\mathbf{u}(\mathbf{r})
\end{align}

即
\begin{equation}
\begin{bmatrix}
\phi_1 \\[4pt]
\phi_2
\end{bmatrix}
=
-
\begin{bmatrix}
k_{1x} & k_{1y} \\[4pt]
k_{2x} & k_{2y}
\end{bmatrix}
\begin{bmatrix}
u_x \\[4pt]
u_y
\end{bmatrix}
\end{equation}


在一级近似下，位移场$\bm{u(r)}$正比于水面高度的斜率，
\begin{equation}
    \bm{u(r)} \approx -H(1-n_a/n_w) \nabla h(\bm{r})
\end{equation}
其中 $h(\bm{r})$ 表示水面高度，$H$ 为水面与背景图案之间的距离，而 $n_a/n_w$ 是空气与水之间的折射率之比。

所以，对$\bm{u(r)}$进行积分，并除以系数，即可得到$h(\bm{r})$，也就是我们要求的水位高度场（相对水位变化）。

% ===========================================
\subsection{直线波源}

为了实现菲涅尔波带片对水面波的聚焦，需要首先获得具有良好横向均匀性的入射波场。理论上，理想的直线波源应满足以下条件：在波源方向上具有近似连续且均匀的振幅分布，在传播方向上产生近似平面波前，同时能够独立控制波的振幅与频率。然而，在实际实验条件下，连续直线型水波源的实现存在一定困难。因此，本实验针对不同波源方案进行了多次改进，以提高入射波场质量，降低非理想因素对波带片聚焦效果的影响。

最初实验采用多个点波源排列的方法近似构造直线波源\cite{Che2024STVP}。具体而言，将8个独立振动点源沿直线方向排列，使各点源同时振动，通过空间叠加形成近似直线传播的水波。该方案基于惠更斯原理，即连续波阵面可以看作由大量次级波源组成，因此有限数量的离散点源在理论上能够近似模拟连续线源。然而，由于实际点源数量有限，点源之间存在明显的空间离散性，其产生的波场并不能完全等效于连续直线波源。

实验观察发现，点源阵列产生的波面存在明显的不连续结构。在相邻点源之间，由于传播距离和相位关系不同，会形成局部干涉增强与减弱区域，导致波前出现周期性的调制结构，即明显的单胞特征。这种离散化误差会破坏入射波的均匀性，使后续波带片受到的激励不再满足理想平面波假设，从而影响聚焦区域的清晰程度。类似的多源干涉现象本质上来源于有限数量相干源之间的空间叠加规律，是水波实验中研究波前重构和干涉现象的重要基础。

为了改善点源阵列造成的离散性问题，随后采用气动线性波源进行实验。该方法通过设计狭长直线型出风口，使空气周期性作用于水面，在气流扰动区域产生连续水面振动，从而形成更加接近理想线源的波场。相比于离散点源，气动波源能够显著提高波源连续性，使产生的波面更加平滑，因此在实验中获得了较好的线性传播效果。

然而，气动波源仍存在若干限制。首先，由于受到出风口尺寸限制，实际气动波源长度小于水池宽度，因此无法完全覆盖整个实验区域。波源两端附近由于边界效应以及气流分布不均，会出现明显的振幅衰减，使传播后的水波在横向方向上存在不均匀性，靠近水池两侧区域的波幅明显降低。这种横向振幅梯度会降低入射波场与理论模型之间的一致性。

其次，气动波源涉及空气--水界面的复杂耦合作用。空气流动产生的剪切应力、压力扰动以及自由液面响应共同决定最终水波振幅，其作用机制受到流速分布、喷口形状以及气流稳定性的共同影响。风生水波的产生过程本身就是一个复杂的两相流动力学问题\cite{Plate1969}，因此相比于直接机械驱动，气动方式难以建立输入参数（例如气流强度）与水波振幅之间简单、稳定的对应关系。此外，实验中气管弯折、气流阻尼以及供气系统波动等因素也容易引入额外误差。

因此，虽然气动波源改善了波源连续性，但仍存在振幅控制困难、空间均匀性不足等问题。

综合考虑波场质量、参数可控性以及实验稳定性，最终采用机械驱动固体线性波源作为实验方案，如图\ref{fig:line_source}所示。该波源由接近水池宽度的三角形截面亚克力棒构成，通过电机驱动其沿竖直方向周期运动，使棒体直接推动水体产生传播波。机械式造波方法是实验水动力学中常用的波浪产生方式，其优点在于可以直接控制边界运动，从而获得稳定、可重复的周期波场\cite{Flick1980}。

\begin{figure}[!h]
    \centering
    \includegraphics[width=0.4\textwidth]{line_source.png}
    \caption{机械驱动固体线性波源} 
    \label{fig:line_source}
\end{figure}

与前两种方案相比，固体线性波源具有更加明确的参数控制关系。首先，在振幅控制方面，实验中测得波源振幅与电机输入振幅基本一致，因此可以直接通过控制电机运动幅度调节水波振幅，避免了气动方式中空气-水耦合带来的不确定性。其次，在频率控制方面，机械驱动频率直接决定波源振动频率，实验有效工作范围为2-6 Hz，生成水波频率与输入控制信号保持一致，能够满足不同波长条件下的实验需求。

此外，由于亚克力棒长度接近水池宽度，其横向尺寸远大于波长尺度，因此产生的水波在横向方向具有较好的均匀性，如图\ref{fig:line}所示。实验观察表明，传播过程中波面的振幅分布较为一致，靠近水池两侧未出现明显衰减，各周期波之间的振幅稳定性也较高。这使得最终获得的入射波更加接近理论模型中的理想直线波，为后续菲涅尔波带片聚焦实验提供了可靠的入射条件。

\begin{figure}[!h]
    \centering
    \includegraphics[width=0.4\textwidth]{line.png}
    \caption{波源造波结果} 
    \label{fig:line}
\end{figure}

综上，本实验通过点源阵列、气动线性波源以及机械固体线性波源三种方案的逐步优化，实现了从离散近似到连续驱动，再到高稳定机械控制的改进过程。最终采用的亚克力棒机械波源兼具振幅、频率可控以及横向均匀性高等优势，能够有效降低非理想入射波场对菲涅尔波带片聚焦实验的影响。

\subsection{三种菲涅尔波带片设计与迭代}

菲涅尔波带片的聚焦效果依赖于各波带对入射波场的调制作用。理想情况下，各波带应满足相位关系，使经过不同路径传播的水波在目标焦点处产生相长干涉。

振幅型波带片通过选择性阻挡部分波传播区域实现空间振幅调制，其工作原理与经典光学振幅型菲涅尔波带片类似，即利用透射区域和遮挡区域之间的差异改变传播波场分布。x 型和 h 型波带片则是在传统振幅调制基础上进一步优化结构形式，通过改变波带排列方式以及局部阻挡区域，使水波传播过程中产生更加稳定的干涉增强效果。三种结构的设计与比较，有助于研究不同空间调制方式对水波聚焦性能的影响。

\subsubsection{振幅型波带片}

根据图\ref{fig:bodaipianshiyitu}的结果，我们可以计算出振幅型波带片每个隔板边缘的坐标$r_m$，满足
\begin{equation}
    \sqrt{{r_m}^2+{R_0}^2}-R_0=m \times \frac{\lambda}{2}
\end{equation}
化简即得
\begin{equation}
    r_m=\sqrt{m\lambda f+(m \cdot \frac{\lambda}{2})^2}
\end{equation}
据此可以设计振幅型波带片。

在这里我们设计了滑槽，可以根据需要便捷地改变滑槽位置，波带片直接插入滑槽固定。

\begin{figure}[htbp]
    \centering
    \begin{minipage}{0.47\linewidth}
        \centering
        \includegraphics[width=\linewidth]{amplitude_plate.png}
        \caption*{(a) 振幅型波带片}
    \end{minipage}
    \hfill
    \begin{minipage}{0.47\linewidth}
        \centering
        \includegraphics[width=\linewidth]{huacao.png}
        \caption*{(b) 滑槽}
\end{minipage}
\caption{振幅型波带片结构图}
\label{fig:amp}
\end{figure}

\begin{figure}[htbp]
    \begin{minipage}{0.47\linewidth}
        \centering
        \includegraphics[width=\linewidth]{x_plate.png}
        \caption*{(a) 相位x型波带片}
    \end{minipage}
    \hfill
    \begin{minipage}{0.47\linewidth}
        \centering
        \includegraphics[width=\linewidth]{h_plate.png}
        \caption*{(b) 相位h型波带片}
    \end{minipage}
    \caption{两种相位型涅尔波带片结构示意图}
    \label{fig:two_fzp}
\end{figure}

\subsubsection{相位型波带片}
相位型波带片通过调节水深以及不同水深的距离来实现相位调节，可以简单分为两种类型，一种是保持水深不变调节通过浅水区的距离，记作相位x型波带片；另一种保持水波通过的距离相同改变水深，记作相位h型波带片.

对相位x型波带片，同样通过求解相位差得到：
\begin{equation}
\Delta x = \frac{2n\pi+fk_1-k_1\sqrt{f^2+y^2}}{k_2-k_1}
\end{equation}
其中，$\Delta x$表示波带片宽度，$f$表示焦距，$k_2$是浅水区波数，$k_1$是深水区波数，$y$是波带片的纵坐标。

对相位h型波带片，计算比较复杂，需要将相位差与色散关系式\eqref{dispersion}联立求解。



\subsubsection{其他说明}

波带片采用 PLA Pro+ 材料进行三维打印加工。聚乳酸（polylactic acid, PLA）是一种常用于熔融沉积成型（Fused Deposition Modeling, FDM）的生物基热塑性聚合物材料，具有加工温度较低、成型精度较高、尺寸稳定性较好等特点，广泛应用于快速原型制造和实验器件制备\cite{Penumakala2020}。本实验选用的 PLA Pro+ 耗材具有较高硬度以及较好的表面平整度，在打印完成后能够保持较稳定的几何尺寸，不易因水流作用发生明显变形，因此适合作为水波实验中的结构材料。

由于实验所使用打印设备的成型尺寸有限，无法一次性打印完整尺寸的波带片。因此，在加工过程中将波带片进行分割打印，并利用滑槽结构和卯榫结构进行拼接组装。为了降低拼接区域对水波传播的影响，所有连接界面均经过机械打磨处理，使拼接面保持平整连续，尽可能减少局部突起、缝隙以及几何不连续性对波场造成的散射和扰动。

此外，由于 PLA 材料密度与水接近，但是打印时往往不是实心的未经处理的波带片容易受到浮力作用产生漂浮或在水波作用下发生移动，从而影响实验稳定性。因此，在波带片内部设计安装铁块作为配重，使其能够稳定放置于水池底部。配重结构不仅提高了器件在实验过程中的稳定性，同时避免了波带片与水面发生非预期相互作用，保证了水波主要受到设计结构本身的调制作用。

综上，本实验通过理论计算、材料选择以及机械结构优化，完成了三种不同类型水波菲涅尔波带片的设计与加工。


\subsection{升级的波带片：正弦平方式相位h型波带片}

在前述三种菲涅尔波带片实验基础上，为进一步提高水波聚焦效果并降低结构边缘引入的非理想扰动，本文对传统波带片结构进行了进一步优化，设计了一种具有两侧缓冲区域的改进型菲涅尔波带片，如图\ref{fig:buffer_fzp}所示。

\begin{figure}[htbp]
    \centering
    \begin{minipage}{0.47\linewidth}
        \centering
        \includegraphics[width=\linewidth]{h_plate.png}
        \caption*{(a) 原始波带片结构}
    \end{minipage}
    \hfill
    \begin{minipage}{0.47\linewidth}
        \centering
        \includegraphics[width=\linewidth]{new_plate.jpg}
        \caption*{(b) 带缓冲区的升级波带片}
    \end{minipage}
    \caption{传统波带片与带缓冲区优化结构对比}
    \label{fig:buffer_fzp}
\end{figure}

传统菲涅尔波带片通常采用明确的几何边界划分不同波带区域，通过阻挡或改变部分区域内水波传播实现空间调制。然而，在实际水波实验中，理想的突变边界会引入额外的边缘散射效应。当水波传播至波带片边缘位置时，由于结构高度和边界条件发生突然变化，局部区域会产生复杂的绕射和反射过程，使部分能量偏离理论设计路径，从而形成额外扰动波场。这些非理想波场可能降低焦点区域的能量集中程度，并影响聚焦结果的稳定性。

针对上述问题，本实验在原有波带片结构基础上，在波带边缘两侧增加缓冲区域，通过平滑化结构过渡降低几何突变带来的散射效应。该设计思想类似于波场调控中的渐变结构，通过减小空间参数变化的不连续程度，使入射水波能够更加平缓地经过调制区域，从而减少高频空间扰动成分。

具体而言，升级后的波带片不再采用完全锐利的边界，而是在原有阻挡区域与开放传播区域之间引入一定宽度的过渡区域，使结构变化更加连续。由于水波传播过程中对边界变化较为敏感，缓冲区能够降低局部反射和散射产生的影响，使更多有效波能量按照设计方向传播，并在目标焦点区域形成增强干涉。

实验结果表明，相比未优化的波带片结构，增加缓冲区域后的波带片能够获得更加明显的聚焦现象。焦点附近的波幅增强更加突出，背景扰动波纹有所减弱，说明结构边缘平滑化能够有效降低非理想散射对水波聚焦过程的影响。

需要指出的是，该优化并未改变菲涅尔波带片基于相位关系的基本聚焦机制，而是在实际实验条件下对理想模型进行了工程修正。由于真实水波系统中存在有限水深、波幅衰减、边界反射以及加工误差等因素，通过引入缓冲区域降低结构引入的额外扰动，有助于提高实验系统与理论模型之间的一致性。

因此，该升级型波带片是在传统菲涅尔波带片基础上的实验优化方案，通过改善结构边界条件提升了水波场调控能力，为后续进一步研究高质量水波聚焦以及基于局域增强波场的粒子操控提供了新的实验基础。


% -------- 四、实验与仿真结果 --------
\section{实验与仿真结果}
我们首先对水波的一些基本参数进行测量，主要是色散关系和衰减系数，是我们的仿真程序使用的参数。仿真程序使用MATLAB编程，由AI辅助完成，具体结构说明将在附录A中介绍。

接着我们对第三部分中设计的四种波带片进行实验分析。水波通过菲涅尔波带片后，由于相位差异会出现干涉图案，肉眼即可观察。为了更清晰的看出水波振幅的变化，我们在实验中拍摄图片，用FCD算法复原出相应水位高度场（相对高度变化），与仿真结果对比，可以清晰看出图案的相似性，在一定程度上验证了理论推导的正确性。

但是图案只是定性描述汇聚现象，我们需要用物理量定量评价聚焦效果。我们仿照光学\cite{hecht2017optics}的定义，并参考Stamnes的推导\cite{Stamnes1986}，使用点扩散函数的相关概念进行说明。此处使用MATLAB程序“focal\_analysis”进行聚焦效果分析，通过对提取出的时间平均的RMS振幅图片进行分析，对焦点进行二维高斯拟合得到焦点坐标、焦点在两个垂直方向的半高宽，如图\ref{fig:Gauss_fit}所示。其中，x坐标对应焦深，x方向的半高宽对应焦深，y方向的半高宽对应焦点的半高宽。

\begin{figure}[!h]
    \centering
    \includegraphics[width=0.4\textwidth]{Gauss_fit.png}
    \caption{高斯拟合} 
    \label{fig:Gauss_fit}
\end{figure}

根据高斯拟合的结果，我们可以对聚焦效率进行计算。仍然是处理时间平均的RMS振幅图片，图片范围为波带片中心后方10cm至40cm左右的范围，宽度约为40cm，左右基本对称。设焦斑面积为$\Omega$，范围为从半高宽到峰值点，设总场面积为$\Omega_0$，定义聚焦效率为
\begin{equation}
\eta=\frac{\iint_{\Omega}[\text{RMS}(x,y)]^2\,dxdy}{\iint_{\Omega_{0}} [\text{RMS}(x, y)]^2 \, dxdy}
\end{equation}
即焦斑面积内的能量与观察范围的能量的比值。聚焦效率越高，说明焦斑能量越集中，聚焦效果越好。

此外，直接比较振幅大小也可以说明聚焦效果。由于实验中焦斑位置几乎不发生改变，且衰减系数随水深变化没有明显的规律性，我们可以直接比较RMS振幅的绝对高度，高度越高，说明能量越大。为了更直观地看出焦斑的变化，我们将随水深变化的结果做成视频，置于PPT中讲解，在这里就只比较聚焦系数的变化。

具体的实验结果以及仿真对比如下：

\subsection{水波衰减系数测量}

\begin{figure}[htbp]
    \centering
    \begin{minipage}{0.47\linewidth}
        \centering
        \includegraphics[width=\linewidth]{decay.png}
        \caption*{(a) 拟合结果}
        \label{fig:decay}
    \end{minipage}
    \hfill
    \begin{minipage}{0.47\linewidth}
        \centering
        \includegraphics[width=\linewidth]{amp_fz.png}
        \caption*{(b) 衰减系数随水深的变化}
        \label{fig:decay_const}
    \end{minipage}
    \caption{衰减系数}
\end{figure}

实际中，由于流体具有阻尼，传播过程中振幅会发生衰减。如图\ref{fig:decay}是直线波源造出的波在水槽中运动的一个剖面振幅图，用带e指数衰减项的振幅公式对其进行拟合，可以得到衰减系数。\ref{fig:decay_const}是对水深从$3.0\texttt{cm}$至$4.0\texttt{cm}$，间隔$0.1\texttt{cm}$的振幅进行拟合，得到的衰减系数散点图，每个点都是五个数据计算的结果。观察到衰减系数随水深无明显变化趋势，在$0.028\texttt{m}^{-1}$附近波动，因此仿真设置的衰减系数为$0.028\texttt{m}^{-1}$。

\subsection{振幅型波带片结果}

\begin{figure}[htbp]
    \centering
    \begin{minipage}{0.47\linewidth}
        \centering
        \includegraphics[width=\linewidth]{amp_ex.png}
        \caption*{(a) 实验结果}
    \end{minipage}
    \hfill
    \begin{minipage}{0.47\linewidth}
        \centering
        \includegraphics[width=\linewidth]{amp_fz.png}
        \caption*{(b) MATLAB仿真结果}
    \end{minipage}
    \caption{振幅型波带片聚焦图案}
    \label{fig:amp_pattern}
\end{figure}

实验中振幅图案呈现收束到中心点再发散的现象，与仿真结果一致。

\begin{figure}[!h]
    \centering
    \includegraphics[width=0.4\textwidth]{amp_P.png}
    \caption{聚焦效率随水深的变化} 
    \label{fig:amp_P}
\end{figure}

振幅型波带片的聚焦效率较低，随水深变化杂乱无章，在$7\%$至$8\%$之间抖动，水深$39\texttt{mm}$时可能为异常值。



\subsection{相位x型波带片结果}

\begin{figure}[htbp]
    \centering
    \begin{minipage}{0.47\linewidth}
        \centering
        \includegraphics[width=\linewidth]{x_ex.png}
        \caption*{(a) 实验结果}
    \end{minipage}
    \hfill
    \begin{minipage}{0.47\linewidth}
        \centering
        \includegraphics[width=\linewidth]{x_fz.png}
        \caption*{(b) MATLABx仿真结果}
    \end{minipage}
    \caption{相位x波带片聚焦图案}
    \label{fig:x_pattern}
\end{figure}

实验中振幅图案呈现中心环状，四周发散的现象，与仿真结果较为相近。

\begin{figure}[!h]
    \centering
    \includegraphics[width=0.4\textwidth]{x_P.png}
    \caption{聚焦效率随水深的变化} 
    \label{fig:x_P}
\end{figure}

相位x型波带片的聚焦效率在设计值上超过了$10\%$,与振幅型波带片相比有接近两倍的提升，并且聚焦效率随水深增加迅速减小至稳定，可能是因为色散关系与水深呈现非线性，当水深过大时，水深变化引起的相位变化迅速减小，以至于没有明显聚焦。与之对比，振幅型波带片虽然最高聚焦效率没有相位x型高，但在较大水深范围内都能出现聚焦现象。

\subsection{相位h型波带片结果}

\begin{figure}[htbp]
    \centering
    \begin{minipage}{0.47\linewidth}
        \centering
        \includegraphics[width=\linewidth]{h_ex.png}
        \caption*{(a) 实验结果}
    \end{minipage}
    \hfill
    \begin{minipage}{0.47\linewidth}
        \centering
        \includegraphics[width=\linewidth]{h_fz.png}
        \caption*{(b) MATLAB仿真结果}
    \end{minipage}
    \caption{相位h波带片聚焦图案}
    \label{fig:h_pattern}
\end{figure}

实验中振幅图案呈现围绕中心的点状分布，与仿真结果相近。

\begin{figure}[!h]
    \centering
    \includegraphics[width=0.4\textwidth]{h_P.png}
    \caption{聚焦效率随水深的变化} 
    \label{fig:h_P}
\end{figure}

相位h型波带片在设计水深的聚焦效率最高，体现了其良好的聚焦能力，且随水深减小的速度明显慢于相位x型波带片，有更强的水深兼容性，但是大范围水深变化下，例如水深达到$40\texttt{mm}$时，效果仍然不及振幅型波带片。

综合来看，相位h型波带片的聚焦效果优于另外两种类型，因而我们在此基础上进行改进，希望通过减小反射来进一步提高聚焦效率。


\newpage

\subsection{正弦平方式相位h型波带片结果}

改进的方法是设计正弦式相位h型波带片，通过边界的平滑过渡来减小反射。

\begin{figure}[htbp]
    \centering
    \begin{minipage}{0.47\linewidth}
        \centering
        \includegraphics[width=\linewidth]{sinh_ex.png}
        \caption*{(a) 实验结果}
    \end{minipage}
    \hfill
    \begin{minipage}{0.47\linewidth}
        \centering
        \includegraphics[width=\linewidth]{sinh_fz.png}
        \caption*{(b) MATLAB仿真结果}
    \end{minipage}
    \caption{正弦平方式相位h型波带片聚焦效果}
    \label{fig:sinh_ex}
\end{figure}

具体实验结果如图\ref{fig:sinh_ex}所示，聚焦图案与相位h型、相位x型波带片的图案相似，但略有不同，呈现从上至下逐渐分叉的结构。

\begin{figure}[!h]
    \centering
    \includegraphics[width=0.4\textwidth]{sinh_P.png}
    \caption{聚焦效率随水深的变化} 
    \label{fig:sinh_P}
\end{figure}

由于正弦式相位h型波带片具有底座，实际水深会高一些，这里的$35\texttt{mm}$水深相当于相位h型波带片的$33\texttt{mm}$水深，但是比较而言，并没有十分显著的聚焦优势。

\newpage

% -------- 五、测量精度与不确定度说明 --------
\section{测量精度与不确定度说明}

在第四部分中我们用到的物理量主要是与点扩散函数相关的焦深、焦距、半高宽，以及与聚焦效率相关的RMS振幅，这说明需要测量的物理量是x、y方向（棋盘格平面）的距离以及水波的相对水位高度，但是在聚焦效率的计算中，影响量是水位高度的相对值（比例大小），因而只需要测量x、y方向的精度与不确定度。

由于实验使用的FCD程序进行了傅里叶变换和积分运算，且实验中还需考虑相机镜头畸变、棋盘格大小、像素误差等的影响，直接计算精度与不确定度非常复杂，因而我们采用比较法进行计算。具体操作是将FCD程序计算得到的值与用直尺、游标卡尺直接测量得到的值进行比较，计算系统误差、合成标准不确定度、扩展不确定度\cite{JJF1001-2011}。

\subsection{测量精度}

保持相机、光源等位置不动，将直尺放在棋盘格上拍照，多次改变直尺位置和FCD程序图片选取范围，根据实验标准偏差计算公式\cite{JJF1001-2011}
\begin{equation}
S = \sqrt{\frac{\Sigma^{n}_{i} (x_i-\bar{x})^2}{n-1}}，\bar{x} = \frac{1}{n}\Sigma^{n}_{i} x_i
\end{equation}
得到$S(x) = 0.255\texttt{cm}$，$S(y) = 0.0259\texttt{cm}$，系统误差分别为$\delta(x)=2.57\%$，$\delta(y)=3.15\%$，精度常由系统误差表达。

\subsection{测量不确定度}

与测量精度时的操作类似，进行多次重复实验，根据A类不确定度公式
\begin{equation}
S = \sqrt{\frac{\Sigma^{n}_{i} (x_i-\bar{x})^2}{n(n-1)}}
\end{equation}
得到$u_A(x) = 0.081\texttt{cm}$，$u_A(y) = 0.0082\texttt{cm}$
由于直尺测量的不确定度为$u_{ruler} = 0.012\texttt{cm}$，所以合成标准不确定度为$u(x) = 0.081\texttt{cm}$，$u(y) = 0.015\texttt{cm}$，对应的扩展不确定度为$U(x) = 0.16\texttt{cm}$，$U(y) = 0.03\texttt{cm}$

测量得到的x、y方向系统误差在3\%左右，且不确定度都小于2mm，而我们实验中考虑的结果都以cm为单位，说明FCD程序的计算对于水波实验是足够的。

\section{误差来源}

\subsection{程序处理误差}

FCD程序中二维傅里叶变换可能存在积分误差，对于平静的水面或倾斜水面误差较大，因而实验前需调平水槽；水槽底部可能发生蠕变，导致各处水深出现微小差异，程序中使用二次曲线拟合尽量减少影响；实际水波焦斑并非完美的高斯型，可能出现拖尾、多峰相连的情况，高斯拟合只是近似，在一定程度上说明聚焦效果；

\subsection{实验误差}
3D打印件的精度、表面粗糙程度可能给水波运动带来影响；水与尺子的吸附可能对水位测量带来影响；海绵的吸波效果影响反射波的大小；波源不一定稳定，需要要考虑一下每个波在相同位置（传播相同时间之后）振幅是否一致，误差多大，用自由波分析。

% -------- 六、总结 --------

\section{总结}
本文将光学中的经典菲涅尔波带片理论引入流体力学体系，通过理论计算、仿真模拟与物理实验验证，成功实现了对水表面波的空间相位调控与能量聚焦。

在实验平台与观测技术的构建上，针对真实流体中易产生的非线性耗散与边界反射问题，本研究参考文献\cite{Wang2024}，设计并搭建了水波实验槽。通过机械驱动的固体线性波源，配合海绵消波与侧壁自由边界优化方案，有效保障了入射波场的平直性与纯净度。在数据采集与处理方面，引入快速棋盘格解调（FCD）技术，实现了对复杂水面高度场的高精度、非接触式测量。误差与不确定度分析表明，该测量系统的相对误差控制在较低范围内，满足宏观水波精细化特征的提取要求。

在波带片的设计与性能评估上，本研究制做了振幅型、相位 $x$ 型、相位 $h$ 型以及正弦平方式相位 $h$ 型四类波带片，实验与仿真的聚焦图案高度吻合，此外，研究还通过二维高斯拟合计算点扩散函数，并对聚焦效率进行了定量的比较分析。具体而言，振幅型波带片虽然绝对聚焦效率相对较低（维持在 7\% 至 8\% 之间），但其物理机制受色散关系非线性的影响最小，因而在较大的水深变化范围内均能保持稳定的聚焦现象。相比之下，相位 $x$ 型波带片通过调控局部相速度，显著提升了峰值聚焦效率（超过 10\%），但其对水深参数极为敏感，水深的微小偏离会导致聚焦效率迅速衰减。相位 $h$ 型波带片则在设计水深下展现出良好的聚焦能力，且相较于相位 $x$ 型，其水深兼容性得到了明显改善，实现了聚焦效率与鲁棒性的平衡，综合性能较好。此外，针对实体台阶造成的边缘散射与反射问题，我们进一步设计了带有流线型缓冲过渡区的正弦平方式相位 $h$ 型波带片。尽管在目前的实验测量范围内，其定量的聚焦效率相较于常规相位 $h$ 型未见显著提升，但从物理图像上看，该设计有效改善了焦点周围背景波场的纯净度。

综上所述，本研究不仅验证了波动干涉理论在复杂水动力学体系中的普适性，也为基于表面波的波前工程提供了器件设计与定量评估的一个方案。相关结论可以为未来高效收集波浪能、设计新型海岸防护结构，以及开发基于局域波场的微粒操纵技术提供参考。


% -------- 七、展望 --------
\section{展望}
尽管本研究在水波菲涅尔波带片的设计与平面波聚焦实验方面取得了一定成果，但真实流体环境与波动物理的复杂性决定了该领域仍有极其广阔的探索空间。基于本课题的现有基础，未来可从实验装置升级、算法逆向设计以及前沿波场探索三方面开展更深入的研究。

在物理实验与硬件装置方面，未来可以尝试搭建高精度的单点或多点的独立波源系统。相比于当前的直线波源，点波源能够产生各向同性的柱面波或球面波前，从而进一步拓展波带片对非平面入射波场的调控能力。此外，波带片材质的微观物理属性对流体边界层的演化也有着不可忽视的作用，后续研究可系统性地探究不同 3D 打印材料（如具备不同表面粗糙度、亲疏水性的树脂与复合材料）对水波反射、界面黏滞耗散及涡流脱落的具体影响，从材料工程层面进一步优化波带片的流体力学性能。

在波带片的结构设计与性能迭代上，未来可引入前沿的机器学习算法进行智能化逆向设计。传统的波带片设计主要依赖于简化的解析公式与经验参数，难以应对复杂的非线性色散与水深梯度的耦合效应。通过结合数值计算软件（如 MATLAB 与高阶流体仿真工具）生成波场演化数据，我们可以训练深度神经网络或采用其他优化算法，自动寻找出最优的连续型微结构表面，从根源上抑制杂散辐射，实现波带片聚焦效率与波场纯净度的提升。

最后，在波场调控的前沿物理现象探索中，可利用改进后的高级波带片产生更为复杂的拓扑水波结构，例如时空涡旋（Spatio-temporal vortices）。通过在波带片的三维构型中引入特定的螺旋形或方位角相位梯度，可赋予局域水波特定的轨道角动量（OAM），进而在水面上激发出具有动态旋转特征的涡旋波束\cite{Che2024STVP}。这不仅能进一步验证宏观流体体系对复杂光场现象的类比能力，还将为基于局域涡旋波场的微粒动态捕获、水上漂浮物的无接触式定向操纵等前沿交叉应用提供全新的物理手段与理论依据。

\newpage

\bibliographystyle{unsrt}
\bibliography{references}

% -------- 八、附录 --------
\appendix
\section{comsol仿真}
\subsection{水波势流理论与comsol造波器仿真}
假设流体是无粘、不可压缩且无旋的。
由于无旋（$\nabla \times \vec{v} = 0$），我们可以引入标量函数速度势$\varphi(x,y,z,t)$，使得流场速度向量 $\vec{v}$ 可以表示为势函数的梯度：

\begin{equation}
\vec{v} = \nabla \varphi = (\frac{\partial \varphi}{\partial x}, \frac{\partial \varphi}{\partial y}, \frac{\partial \varphi}{\partial z})
\end{equation}

由不可压缩条件（$\nabla \cdot \vec{v} = 0$），我们可以得到整个流场内部必须满足的主控方程——拉普拉斯方程\cite{Liang2009}：

\begin{equation}
\nabla^2 \varphi = 0
\end{equation}

水槽的底部和侧壁是不透水的刚性壁面，流体微团不能穿透壁面，也不能脱离壁面，所以在壁面上，流体的法向速度必须等于壁面的法向速度。

由于水底和侧壁是静止的，其速度为 0。设壁面的单位法向量为 $\vec{n}$，则流体在壁面处的法向速度为 0：

\begin{equation}
\vec{v} \cdot \vec{n} = 0
\end{equation}

代入速度势的定义 $\vec{v} = \nabla \varphi$，得到：

\begin{equation}
\nabla \varphi \cdot \vec{n} = 0 \quad \Rightarrow \quad \frac{\partial \varphi}{\partial n} = 0
\end{equation}

对应COMSOL中“零通量”边界条件。

造波器表面的法向速度为$V_0 \cos(\omega t)\sin(\frac{\alpha}{2})$，流体与壁面法向速度相等，因此边界条件为：
\begin{equation}
    \frac{\partial \varphi}{\partial n} = V_0 \cos(\omega t)\sin(\frac{\alpha}{2})
\end{equation}
其中$\alpha$为三棱柱的顶角。

液体自由表面的边界条件为：
\begin{equation}
    \left. \frac{\partial \varphi}{\partial z} \right|_{z=\xi} = \frac{\partial \xi}{\partial t} + \left. \frac{\partial \varphi}{\partial x} \right|_{z=\xi} \frac{\partial \xi}{\partial x} + \left. \frac{\partial \varphi}{\partial y} \right|_{z=\xi} \frac{\partial \xi}{\partial y}
\end{equation}
\begin{equation}
    \left( \frac{\partial \varphi}{\partial t} + \frac{1}{2} \left[ \left( \frac{\partial \varphi}{\partial x} \right)^2 + \left( \frac{\partial \varphi}{\partial y} \right)^2 + \left( \frac{\partial \varphi}{\partial z} \right)^{2} \right] \right) \Bigg|_{z=\xi} + g\xi = 0
\end{equation}
其中(12)式为运动学条件，(13)式为动力学条件\cite{Liang2009}。

对于微幅波，水面位移 $\xi(x,y,t)$ 很小，速度 $v_x=\left. \frac{\partial \varphi}{\partial x} \right|_{z=\xi},v_y=\left. \frac{\partial \varphi}{\partial y} \right|_{z=\xi}$ 也很小。因此速度与波高梯度的乘积项是二阶小量，可以略去，同时$\left. \frac{\partial \varphi}{\partial z} \right|_{z=\xi}$可近似为$\left. \frac{\partial \varphi}{\partial z} \right|_{z=0}$：
\begin{equation}
    \left. \frac{\partial \varphi}{\partial z} \right|_{z=\xi} = \frac{\partial \xi}{\partial t}
\end{equation}
速度平方项$\frac{1}{2} \left[ \left( \frac{\partial \varphi}{\partial x} \right)^2 + \left( \frac{\partial \varphi}{\partial y} \right)^2 + \left( \frac{\partial \varphi}{\partial z} \right)^{2}\right]$也是二阶小量，可以略去，同时$\left. \frac{\partial \varphi}{\partial t} \right|_{z=\xi}$可近似为$\left. \frac{\partial \varphi}{\partial t} \right|_{z=0}$：
\begin{equation}
    \left. \frac{\partial \varphi}{\partial t} \right|_{z=0} + g\xi = 0 
\end{equation}
联立略去二阶小量后的边界条件，可得综合自由液面边界条件：
\begin{equation}
    \left. \frac{\partial \varphi}{\partial z} \right|_{z=0} = -\frac{1}{g}\left.\frac{\partial^2 \varphi}{\partial t^2}\right|_{z=0}
\end{equation}

我们要转入频域求解，将造波器的简谐运动用复数指数形式表示：
\begin{equation}
    V_z(t) = V_0 \cos(\omega t) = \text{Re}\{V_0 e^{i\omega t}\}
\end{equation}
同样地，将时域的速度势 $\varphi(x,y,z,t)$ 分离为空间复振幅 $\phi(x,y,z)$ 与时间项 $e^{i\omega t}$ 的乘积：
\begin{equation}
    \varphi(x,y,z,t) = \text{Re}\{\phi(x,y,z) e^{i\omega t}\}
\end{equation}
将其代入边界条件方程(11)(16)中，得到频域的边界条件方程：
\begin{equation}
    \frac{\partial \phi}{\partial n} = V_0\sin(\frac{\alpha}{2})
\end{equation}
\begin{equation}
    \left. \frac{\partial \phi}{\partial z} \right|_{z=0} = \frac{\omega^2}{g} \phi|_{z=0}
\end{equation}

求解出势函数的复振幅 $\varphi$ 后，根据式()就可以求解出水面的位移。瞬时水面位移为
\begin{equation}
    \xi = -\frac{1}{g}\left.\frac{\partial \varphi}{\partial t}\right|_{z=0}
\end{equation}
代入$\varphi = \text{Re}\{ \phi e^{i\omega t} \}$可得：
\begin{equation}
    \xi = \frac{\omega}{g} \text{Im}\{ \phi e^{i\omega t} \}
\end{equation}
所以振幅就是：
\begin{equation}
    \vert{}A\vert{} = \frac{\omega}{g} \vert{}\phi\vert{}
\end{equation}



\subsection{comsol仿真结果}
根据前文的水波势流理论，可以使用comsol二维-数学-经典偏微分方程-拉普拉斯方程进行造波器的仿真。
忽略流体的粘性，进行几何建模并设置边界条件后即可求解速度势和水面振幅。构造驻波场，以便于观察水面波纹是否是简谐波，而不会产生非线性现象。
如图4，设置区域的长度为5个波长，其余设置与实验保持一致。
\begin{figure}
    \centering
    \includegraphics[width=0.75\linewidth]{fresnel0719-1.png}
    \caption{水面振幅仿真结果}
    \label{fig:placeholder1}
\end{figure}

求解出的水面振幅如图5所示，可以看出其是较好的简谐波，验证了造波器的可行性。
\begin{figure}
    \centering
    \includegraphics[width=0.75\linewidth]{fresnel0719.png}
    \caption{水面振幅仿真结果}
    \label{fig:placeholder}
\end{figure}


\section{MATLAB仿真程序说明}
根据波带片的种类，这里分为了振幅型波带片仿真control\_amp\_new、相位x型波带片仿真control\_phasex\_new、相位h型波带片control\_phaseh\_new仿真三类，均为AI辅助的自写代码。程序是为了观察在水深改变的情况下，波带片聚焦效果的变化。

三个代码的结构几乎一致，只有在波带片设计上依据3.3有变化，因此就一起介绍。

首先是物理参数设置，根据查表获得的以及实验得到的值进行设置。

接着对水面进行网格化处理，依据3.3中的原理按照输入的设计参数设计三种类型的波带片。然后利用for循环不断改变水深，在每一个循环内部，计算出相应的水深矩阵、相速度矩阵，通过FDTD模拟，直接在时间域和空间域中数值求解拉普拉斯方程，计算得到每个时间步的水位矩阵U3。

最后保存结果，一方面是RMS振幅（平方后即为能量），保存的三个数据分别对应整个水面的RMS振幅、波带片后方的RMS振幅、接近实验观察范围的RMS振幅数据，方便进行聚焦效果的分析；另一方面是最后一帧的振幅，包括整个水面的振幅、波带片后方的振幅，方便与实验照片处理得到的图案进行对比。


\section{MATLAB测量程序说明}
程序参考了Wildeman, Sander关于快速棋盘格解调的文章，原理已在3.1中介绍过，这里主要介绍程序
\subsection{调用的函数}
文件夹functions里面是测量程序调用的函数。其中函数fcd\_dispfield3、fcd\_phasefield2、fftinvgrad2、findorthcarrierpks、findpeaks2、getcarrier2、kvec、phase2disp2、subplane、unwrap2是已有的函数，其余的函数，即extractline\_function、extractline\_function2、extractpoint\_function、extractpoint\_function2、format\_spatiotemporal\_axes、plot\_wave\_field、points\_function、save\_fig\_png为自写函数，使用AI辅助编程，具体过程是根据代码需要构思函数功能与结构，将想法告诉AI，AI生成初版代码，我们阅读代码并检查是否有问题，运行代码，根据运行结果考虑进一步优化代码，手动修正与AI编程同步进行。
\subsection{主程序}

find\_points程序是一个提取坐标像素点的程序，调用find\_points函数，用于划定照片裁剪范围，也就是实验观察范围。

water\_level\_restoration\_new程序是一个用于处理照片的程序，原理参照3.1，代码结构分为5部分。

第一部分是设置。首先输入保存的实验图片的路径以及名称等信息，便于后续提取处理；接着设置输出信息，可根据需要打开开关，设置想要的输出信息；最后是输入参数，是与图片处理相关的物理参数。

第二部分是图像处理。首先处理静态图片（每次可处理一张静态图片与多张动态图片），先利用points\_function在静态图片上划分观察范围，接着裁剪掉边缘部分防止带来后续计算问题，然后用findorthcarrierpks函数从静态图像中提取两个正交方向，再用getcarrier2函数提取这两个方向的倒易晶格向量。动态图像由于有水面的变化，存在一个位移场$u(r) = (u, v)$，相当于是静态图像的失真，利用调制解调的原理，经过 fcd\_dispfield3、fcd\_phasefield2、phase2disp2 函数可求得位移场$u(r) = (u, v)$，最后通过函数 fftinvgrad2 进行二维傅里叶积分计算出水位场 h。

第三部分是绘图。首先确定保存图片的信息等，主要来源于第一部分的设置。第二部分求出的h是像素单位，绘图前还需要进行转化。运用 subplane 函数去除掉平面，并利用最小二乘法拟合水槽底面的弯曲，再乘上单位转化系数 mm\_per\_pixel 得到实际水面高度变化量（水面波动）height，并进行绘图。

第四部分是数据提取与多帧保存。利用处理的第一帧动态图片进行操作，可以选择提取水面上的点或者线的位置、高度信息，后续处理的动态图像会一起保存在一个文件中。

第五部分是数据保存。将第四部分提取的数据写入CSV文件中，方便后续进行数据处理。

sequence\_restoration\_new程序与water\_level\_restoration\_new原理是一致的，在此处就不重复介绍了，仅仅说明其中不同的地方。

sequence\_restoration\_new程序是用于批处理的，将时间上连续的图片按顺序依次处理，得到的位移场信息、水位高度场信息以时间为第三个维度存储在三维矩阵中进行运算，最后再依次提取出成我们想要的信息，如水位高度场、点和线的数据，优势在于处理速度更快。

为了方便后续分析聚焦效率，加入了RMS能量分析代码，也就是均方根振幅求解，具体是将水位高度场的结果平方，逐帧相加，最后除以帧数得到，通过save\_fig\_png函数保存成一个二维图片和.fig文件，二维图片可用于后续合成视频分析，.fig文件用于聚焦效率分析。

\section{MATLAB聚焦效率程序}
\subsection{调用的函数}
聚焦效率分析所用的两个函数calc\_water\_wave\_efficiency、Guass\_peak\_analysis\_fig均为AI辅助的自写代码。

Guass\_peak\_analysis\_fig函数用于对焦斑进行二维高斯拟合。使用者选择.fig文件，自行点击焦斑，程序会使用二维高斯函数自动拟合，并输出水波的焦点坐标、半高宽、拟合振幅及其误差等的信息，并保存在相应的二维图片、三维.fig文件。这里面调用了三个函数，辅助函数黑塞矩阵通过二阶偏导数计算误差，辅助函数掩膜加权罚函数用于阻止拟合结果超出物理实际，通过加权实现。

调用的第三个函数是clac\_water\_wave\_efficiency函数，它会接收Guass\_peak\_analysis\_fig函数的拟合结果，利用半高宽和焦点峰值坐标确定焦斑面积，对其积分得到焦斑能量，再除以全场总能量得到聚焦效率，并输出相关参数，返回到Guass\_peak\_analysis\_fig函数。因而Guass\_peak\_analysis\_fig函数实际输出的值还包括了聚焦效率等参数。这里调用了两次lac\_water\_wave\_efficiency函数，是因为对于焦斑能量由不同的定义以及分析，可以使用半高宽计算焦斑面积进而算能量，也可以使用1.5倍半高宽计算，前者考虑能量核心区的聚焦效果，后者考虑整个主峰的聚焦效果。

\subsection{主程序}
focal\_analysis是方便做实验的同学使用的数据分析程序，把需要修改的地方写在最前面，能够让同学便捷地操作。程序的目的是将Guass\_peak\_analysis\_fig函数的分析结果按照水深的递增保存在同一个CSV文件中，便于后续数据处理与分析。

首先是保存文件的路径设置，接着运用for循环，将Guass\_peak\_analysis\_fig函数的分析结果依次写入矩阵中，最后把矩阵转换成表格，保存成CSV文件。



\end{document}

